% =====================================================================
%  GLOSARIO Y ABREVIATURAS (versión de 60 páginas)
%  Solo los términos que hacen falta para leer el resumen y la introducción;
%  los demás se definen en el marco conceptual.
% =====================================================================
\ustaseccion{Glosario}

\setstretch{1.1}
\begin{description}[style=unboxed,leftmargin=0pt,itemsep=0em,parsep=0pt]
  \item[Agente.] Cada uno de los dos robots móviles del sistema, construidos sobre AWS DeepRacer y dedicados cada uno a un piso.
  \item[Cinemática tipo automóvil (dirección Ackermann).] Modelo de un vehículo que dirige con el eje delantero, tiene radio de giro mínimo y no puede girar sobre su propio eje.
  \item[Continuidad del servicio.] Propiedad de no interrumpir la asistencia mientras el usuario cambia de piso.
  \item[Localización.] Estimación de la posición y la orientación del agente respecto a un marco de referencia fijo.
  \item[Misión.] Una solicitud de guiado completa, desde que el usuario elige origen y destino hasta que el sistema la cierra.
  \item[Odometría.] Estimación del desplazamiento del agente a partir de sus propias mediciones; aquí, de barridos sucesivos del láser.
  \item[Punto de transferencia.] Lugar físico, en este trabajo la escalera, donde un agente entrega el guiado y el otro lo recibe.
  \item[Relevo.] Traspaso de la responsabilidad de guiar a un usuario, de un agente al otro, en un punto acordado de antemano.
  \item[Tasa de éxito.] Proporción de misiones exitosas sobre las ejecuciones válidas.
  \item[Tiempo de respuesta.] Intervalo entre la solicitud del usuario y el instante en que el agente asignado arranca la tarea.
  \item[Usuario.] Persona que pide el servicio desde la interfaz y a la que los agentes guían hasta su destino.
  \item[Zona de transición.] Punto del edificio donde se cambia de piso. En este sistema ningún agente la atraviesa.
\end{description}

\ustaseccioncontinua{Abreviaturas}

\setstretch{1.1}
\begin{description}[style=unboxed,leftmargin=0pt,itemsep=0em,parsep=0pt]
  \item[AMCL.] \textit{Adaptive Monte Carlo Localization}, localización adaptativa por método de Monte Carlo.
  \item[AWS.] \textit{Amazon Web Services}, fabricante de la plataforma DeepRacer.
  \item[DDS.] \textit{Data Distribution Service}, protocolo de transporte sobre el que ROS~2 implementa sus comunicaciones.
  \item[GED.] Grupo de Estudio y Desarrollo en Robótica.
  \item[HRI.] \textit{Human-Robot Interaction}, interacción humano-robot.
  \item[IMU.] \textit{Inertial Measurement Unit}, unidad de medición inercial.
  \item[JSON.] \textit{JavaScript Object Notation}, formato de intercambio de datos en texto.
  \item[Nav2.] \textit{Navigation2}, marco de navegación autónoma de ROS~2.
  \item[ROS.] \textit{Robot Operating System}. La versión empleada es ROS~2.
  \item[SLAM.] \textit{Simultaneous Localization and Mapping}, localización y mapeo simultáneos.
\end{description}
\onehalfspacing
